\documentclass[12pt]{article}
\usepackage{geometry}
\geometry{a4paper, margin=1in}
\usepackage{hyperref}
\usepackage{booktabs}
\usepackage{amsmath}

\title{Thesis Scratchpad: Experiment Design for FP8 Inference Optimization of GenomeOcean-4B}
\author{Research Notes}
\date{\today}

\begin{document}
\maketitle

\tableofcontents
\vspace{2em}

%=============================================================================
\section{Overview: The Inference Optimization Hierarchy}
%=============================================================================

This document serves as the technical narrative and experiment log for the FP8 inference optimization experiments conducted on the GenomeOcean-4B Genomic Foundation Model (GFM) using NVIDIA H100 80GB hardware. It is organized as an introduction to the experiment design and is intended for inclusion in the thesis methodology chapter.

The central research question is: \textit{How far can we push the inference throughput of a large genomic foundation model through progressive precision reduction, without compromising its underlying biological representational fidelity?}

To answer this question, we define a progressive hierarchy of three quantization protocols, organized from the most conservative (maximum fidelity) to the most aggressive (maximum throughput):

\begin{table}[h]
\centering
\begin{tabular}{llll}
\toprule
\textbf{Protocol} & \textbf{Name} & \textbf{Weights} & \textbf{KV Cache / Activations} \\
\midrule
B.1 & KV-Cache Quantization (Storage-First)  & BFloat16 (W16) & FP8 (A8) \\
B.2 & W8A8 Dynamic Quantization (Compute-First) & FP8 (W8) & FP8-Dynamic (A8) \\
B.3 & Full FP8 (Zero 16-bit Components)     & FP8 (W8) & FP8 (A8), lm\_head: FP8 \\
\bottomrule
\end{tabular}
\caption{Overview of the three inference optimization protocols evaluated in this study.}
\end{table}

Additionally, \textbf{Protocol A} refers to the PyTorch-native preprocessing experiments (embedding extraction and quality evaluation at full BFloat16 precision) which establish the baseline biological fidelity metrics. These are discussed in the preceding experimental chapter.


%=============================================================================
\section{Protocol B.1: KV-Cache Quantization (Storage-First)}
%=============================================================================

\subsection{Intuition and Motivation}

During autoregressive genomic sequence generation, the bottleneck is not raw compute but \textbf{memory bandwidth}. Each newly generated token requires the GPU to reload the entire accumulated Key-Value (KV) cache—which encodes all prior biological context—from High-Bandwidth Memory (HBM). For a 4-billion parameter model tasked with generating long microbial genome sequences (up to 10,240 tokens $\approx$ 51 kilobase pairs), this KV cache can occupy tens of gigabytes.

The core insight of Protocol B.1 is that the \emph{storage format} of the KV cache can be decoupled from the \emph{compute precision} of the attention mechanism. By instructing vLLM to store the attention keys and values in FP8-E4M3 format at inference time, we compress the KV cache from 16 bits to 8 bits per element---halving the data volume without modifying the model weights at all. The attention dot-product math is still performed in BFloat16 (the 8-bit tensors are dequantized on-the-fly into registers), making this a \textbf{``Storage-First''} optimization: we save on memory bandwidth without using 8-bit tensor cores for the actual computation.

This is intentionally the most conservative protocol. Since the model weights remain in BFloat16 and the compute remains in full precision, the risk to biological fidelity is minimal and easily measurable.

\subsection{Why This Protocol?}

Protocol B.1 establishes the \textbf{baseline for what is achievable with zero model modification}. It answers the question: ``If we change nothing about the model itself, what throughput gain is possible simply by compressing the memory traffic?''

This is the pragmatically safest starting point for a research context where biological fidelity must be verifiable and where the cost of a mistake (e.g., generating taxonomically incorrect sequences) would undermine the entire evaluation pipeline.

\subsection{Benchmark Design}

The benchmark is structured in two phases:

\begin{enumerate}
    \item \textbf{Phase A -- Biological Fidelity (Quality)}: A sliding-window perplexity calculation is performed over 96 randomly sampled microbial genomes ($N=96$). Each genome is tokenized and chunked using a context window of 10,240 tokens ($\approx$ 51 kilobase pairs) with a stride of 2,560 tokens. The NLL (Negative Log-Likelihood) is averaged across all windows per genome. This measures whether the model's prediction distribution has changed with FP8 KV storage.
    \item \textbf{Phase B -- Efficiency (Throughput)}: A saturated single-request throughput benchmark. The $N = 96$ genomes are each given a 10,240-token prefix, and the model is asked to generate 2,560 tokens per genome. The full workload of 96 sequences (at maximum concurrency, batch size 96) is submitted in one call ($N \times K = 96 \times 5 = 480$ total sequences with the 5x workload multiplier). The total generation time divided by total tokens generated yields the honest steady-state throughput.
\end{enumerate}

Key implementation choices that promote measurement integrity:
\begin{itemize}
    \item \textbf{Prefix caching disabled}: \texttt{enable\_prefix\_caching=False} ensures the engine cannot re-use previous KV computations between repeats, forcing genuine memory bandwidth utilization.
    \item \textbf{Warmup separation}: A single-sequence generation run is executed before the timed call. This burns in the CUDA graphs and JIT-compiled kernels so they do not appear in the measured latency.
    \item \textbf{Saturated workload}: The maximum batch size for each precision mode is used (B.1: 96, BF16 baseline: 48), ensuring the H100 HBM and compute units are fully saturated.
\end{itemize}

\subsection{Partial Results}

\begin{table}[h]
\centering
\begin{tabular}{lcc}
\toprule
\textbf{Metric} & \textbf{BFloat16 Baseline} & \textbf{Protocol B.1 (FP8 KV)} \\
\midrule
Max Batch Size             & 48    & 96    \\
Throughput (tok/sec)       & 1829.18 & 2042.66 \\
Energy Efficiency (tok/W)  & 2.64  & 3.14  \\
VRAM Pool Allocated        & 65.83 GB & 65.83 GB \\
\midrule
Perplexity (PPL) $\pm$ std    & 127.68 $\pm$ 52.05 & 128.33 $\pm$ 52.14 \\
NLL $\pm$ std                 & 4.76 $\pm$ 0.42 & 4.77 $\pm$ 0.42 \\
\bottomrule
\end{tabular}
\caption{Protocol B.1 vs. BFloat16 Baseline. Both modes use the \emph{identical model weights}; the only change is the KV cache storage format. The PPL difference is $<$0.52\%, indicating statistically negligible fidelity loss.}
\end{table}

Notably, the \textbf{maximum batch size doubles from 48 to 96} under B.1 (with the VRAM pool remaining constant at 65.83 GB), because each KV entry consumes half the memory. This effectively doubles the biological context capacity per H100, which is a key result for practical discovery pipeline deployment.


%=============================================================================
\section{Protocol B.2: W8A8 Dynamic Quantization (Compute-First)}
%=============================================================================

\subsection{Intuition and Motivation}

Protocol B.1 achieves a ``Storage Victory'': it saves bandwidth but does not engage the H100's specialized FP8 tensor cores (the SM90 accelerators), which are capable of performing fused 8-bit matrix multiplications at up to 989 TFLOPS---approximately 2x the throughput of BFloat16 operations. To access this hardware capability, the model's \textbf{weights themselves} must be quantized to FP8, not just the KV cache.

Protocol B.2 applies \textbf{W8A8 (Weight 8-bit, Activation 8-bit) Post-Training Quantization} via the \texttt{llm-compressor} library. This uses the \texttt{FP8\_DYNAMIC} scheme:
\begin{itemize}
    \item \textbf{Weights}: Quantized \emph{statically} on a per-channel basis using FP8-E4M3. The scaling factors are computed once from the weight tensors and saved into the model checkpoint. This is a one-time offline operation.
    \item \textbf{Activations}: Quantized \emph{dynamically} during inference on a per-token basis. The scaling factor for each activation vector is recomputed from its observed min/max at runtime, ensuring accuracy for unseen biological sequences.
\end{itemize}

\subsection{Why Dynamic Activation Quantization? (The Calibration Problem)}

Static activation quantization requires a \textbf{calibration dataset}: a small, representative set of samples fed through the model to observe the real-world distribution of activation magnitudes. The per-layer scaling factors are then fixed based on these observed ranges.

For a Genomic Foundation Model trained on a vast and taxonomically diverse metagenomic corpus, composing such a calibration dataset is non-trivial:
\begin{itemize}
    \item \textbf{Taxonomic coverage}: A calibration set that overrepresents common phyla (e.g., Proteobacteria) would produce poorly calibrated scaling factors for rare microbial taxa, causing outsized quantization error on precisely the biological sequences the model must generalize to.
    \item \textbf{Sequence diversity}: Unlike natural language, genomic sequences exhibit extreme compositional heterogeneity (GC content, codon bias, repeat structures) across different organisms. A small calibration set cannot capture this diversity.
\end{itemize}

By opting for \textbf{dynamic scaling}, we eliminate the calibration problem entirely. Each forward pass re-evaluates its own activation range based on the actual genomic input, making the system robust to any unseen species or sequence composition. This makes Protocol B.2 a \textbf{``data-free''} quantization approach.

\subsection{Why Exclude the \texttt{lm\_head}?}

The \texttt{lm\_head} is the final linear projection layer that maps the model's internal hidden state (a 3072-dimensional vector) to the vocabulary logit space (the probability distribution over the next biological token, A, C, G, T, etc.). This layer directly determines the biological ``grammar'' of generated sequences.

Because the vocabulary of a GFM encodes subtle evolutionary probabilities (e.g., the probability of a specific codon in a highly conserved gene region), quantizing the \texttt{lm\_head} to 8-bit precision introduces rounding noise into the logit space that could bias generation towards more ``common'' sequences and away from the rare but biologically meaningful ones.

Since the \texttt{lm\_head} accounts for less than 1\% of the model's total parameter count, keeping it in BFloat16 incurs essentially zero performance penalty while preserving the high dynamic range required for accurate biological synthesis.

\subsection{Benchmark Design}

The Phase A and Phase B methodology is identical to Protocol B.1. The key difference is the model loaded by the vLLM engine:
\begin{itemize}
    \item \textbf{B.1 Engine}: Loads \texttt{DOEJGI/GenomeOcean-4B} (raw BFloat16 weights), with \texttt{kv\_cache\_dtype="fp8"}.
    \item \textbf{B.2 Engine}: Loads a locally quantized checkpoint (\texttt{GenomeOcean-4B-FP8-W8A8}), with \texttt{kv\_cache\_dtype="fp8"}. vLLM detects the \texttt{"quantization": "fp8"} field in the checkpoint's \texttt{config.json} and automatically routes matrix multiplications through FP8 GEMM kernels.
\end{itemize}

The quantization is performed offline using:
\begin{verbatim}
python benchmark/quantize_genomeocean_fp8.py --model "DOEJGI/GenomeOcean-4B"
\end{verbatim}
This script applies the \texttt{FP8\_DYNAMIC} recipe via \texttt{llm-compressor}, targeting all \texttt{Linear} layers except the \texttt{lm\_head}.

\subsection{Partial Results}

\textit{[Results pending full saturated HPC run. Benchmark execution command:]}

\begin{verbatim}
python run_generation_benchmark.py \
    --csv ".../bac120_100_seq_550kbp_generation_use.csv" \
    --model "DOEJGI/GenomeOcean-4B" \
    --model-w8a8 "./GenomeOcean-4B-FP8-W8A8" \
    --quant-modes fp8_v2 \
    --batch-sizes 96 \
    --n-genomes 96 --tokens-per-genome 102400 \
    --context-len 10240 --stride 2560 --gen-len 2560 \
    --precision bfloat16 --n-repeats 5
\end{verbatim}

Expected outcomes based on H100 FP8 kernel benchmarks:
\begin{itemize}
    \item \textbf{Throughput}: Expected improvement over B.1 due to reduced weight-loading overhead and FP8 GEMM kernel engagement.
    \item \textbf{Biological Fidelity}: Expected PPL degradation of $<$1\% relative to BF16 baseline (per prior W8A8 literature on language models).
\end{itemize}


%=============================================================================
\section{Protocol B.3: Full FP8 (Zero 16-bit Components)}
%=============================================================================

\subsection{Intuition and Motivation}

Protocols B.1 and B.2 both retain at least one BFloat16 component: B.1 preserves the weights, and B.2 preserves the \texttt{lm\_head}. Protocol B.3 asks the natural maximal question: \textit{What happens to efficiency and biological fidelity if we eliminate all 16-bit components entirely?}

This protocol is not about achieving the ``best'' biological fidelity---by definition it will be worse than B.2. It serves as an \textbf{ablation study}: by measuring the marginal fidelity cost of quantizing the \texttt{lm\_head}, we can quantify how much biological precision is traded away for maximum throughput. This provides a principled bound on the ``compute-first'' optimization frontier.

\subsection{Design}

Protocol B.3 is identical to B.2 with one change: the \texttt{lm\_head} is \textbf{included} in the quantization recipe. This requires a separate quantized checkpoint where the \texttt{ignore} list in \texttt{llm-compressor} is removed:

\begin{verbatim}
recipe = QuantizationModifier(
    targets="Linear",
    scheme="FP8_DYNAMIC",
    ignore=[]  # <-- lm_head is now included
)
\end{verbatim}

\subsection{Expected Observations}

\begin{itemize}
    \item \textbf{Memory}: Negligible additional savings ($<$50 MB) since the \texttt{lm\_head} is a small layer.
    \item \textbf{Throughput}: Minimal additional gain for the same reason.
    \item \textbf{Biological Fidelity (PPL)}: The critical measurement. If the PPL degradation relative to B.2 is small (e.g., $<$1\%), it suggests that the \texttt{lm\_head} is robust to quantization and B.3 is a safe configuration. If the degradation is large, it validates the B.2 design choice to preserve the ``biological output head'' in full precision.
\end{itemize}

\subsection{Partial Results}

\textit{[Pending experiment execution.]}


%=============================================================================
\section{Appendix: Benchmark Infrastructure Summary}
%=============================================================================

\subsection{Hardware and Software Stack}

\begin{itemize}
    \item \textbf{GPU}: NVIDIA H100 80GB SXM (SM90 architecture, HBM3, 3.35 TB/s bandwidth)
    \item \textbf{Inference Engine}: vLLM (latest, with \texttt{VLLM\_ATTENTION\_BACKEND=FLASH\_ATTN})
    \item \textbf{Quantization Tool}: \texttt{llm-compressor}
    \item \textbf{Attention Backend}: FlashAttention-3 (FA3) for FP8 KV layout support
    \item \textbf{Model}: DOEJGI/GenomeOcean-4B (Transformer Decoder, 4B params, GQA with 8 KV heads, 40 attention heads, 40 layers, hidden dim 3072)
\end{itemize}

\subsection{Why H100 is Required}

The H100 is architecturally necessary---not merely preferred---for these experiments:
\begin{enumerate}
    \item \textbf{HBM3 Bandwidth (3.35 TB/s)}: The memory subsystem must be fast enough to stream the multi-hundred-GB KV cache required by Protocol B's batch sizes. Lower-bandwidth GPUs (e.g., A100 at 2 TB/s) would further exacerbate the bandwidth bottleneck being measured.
    \item \textbf{SM90 FP8 Tensor Core Support}: Only SM90 (Hopper) and Ada Lovelace GPUs have native 8-bit GEMM hardware. Protocols B.2 and B.3 are designed to engage these cores. On A100, the FP8 GEMM kernels would fall back to emulation.
    \item \textbf{PagedAttention + FP8 KV Descriptor Support}: The vLLM/FlashAttention-3 FP8 memory layout requires hardware-level support for 8-bit cache descriptors in the memory hierarchy, which is native only to H100.
\end{enumerate}

\end{document}
